\documentclass[11pt]{amsart}
\usepackage{amsmath,amssymb}
\setlength{\emergencystretch}{3em}
\tolerance=600
\usepackage{tikz-cd}
\usepackage{tikz}
\usepackage[hidelinks]{hyperref}

\newtheorem{thm}{Theorem}[section]
\newtheorem{lem}[thm]{Lemma}
\newtheorem{cor}[thm]{Corollary}
\newtheorem{prop}[thm]{Proposition}
\theoremstyle{definition}
\newtheorem{defi}[thm]{Definition}
\newtheorem{rem}[thm]{Remark}
\newtheorem{ex}[thm]{Example}
\newtheorem{que}[thm]{Question}

\newcommand{\Set}{\mathbf{Set}}
\newcommand{\Fin}{\mathbf{FinSet}}
\newcommand{\Sub}{\mathrm{Sub}}
\newcommand{\Aut}{\mathrm{Aut}}
\newcommand{\PP}{\mathsf P}
\newcommand{\HH}{\mathsf H}
\newcommand{\cX}{\mathcal X}
\newcommand{\cC}{\mathcal C}
\newcommand{\cP}{\mathcal P}
\newcommand{\cS}{\mathcal S}
\newcommand{\cK}{\mathcal K}
\newcommand{\cJ}{\mathcal J}
\newcommand{\LC}{\mathbf{LC}}
\newcommand{\id}{\mathrm{id}}


\begin{document}

\title[Intermediate logics and theories with a finite category of models]{Intermediate
logics and the imaginaries of a theory with a finite category of
models: a classification}
\author{Ryota Shibusawa}
\thanks{Corresponding author. E-mail: \texttt{r-shibusawa@daiichi-koudai.ac.jp}}
\address{Daiichi Institute of Technology}
\email{r-shibusawa@daiichi-koudai.ac.jp}
\keywords{intermediate logic, weak excluded middle, G\"odel--Dummett logic, Heyting category, pretopos completion, imaginaries, Kripke model, EI-category}
\subjclass[2020]{03G30 (primary), 03B55, 18F10, 18B25 (secondary)}

\begin{abstract}
For an intermediate logic $L$ and a coherent theory, reading the
theory in first-order logic over $L$ and adding the imaginaries are two
operations which need not commute. We determine when they commute for
the theories whose category of models is a finite category $\cX$ in
which idempotents split, the theory being that of the finite
set-valued functors on $\cX$. For classical logic they always commute.
If $L$ does not prove the weak excluded middle, they commute exactly
when $\cX$ is a groupoid. If $L$ proves the weak excluded middle and
not the linearity axiom, they commute exactly when every endomorphism
of $\cX$ is invertible and no two non-invertible arrows are composable.
If $L$ proves linearity and is not classical, they commute exactly
when every endomorphism is invertible and a combinatorial condition on
the strings of non-invertible arrows with at most $n_L$ objects holds,
which is empty for the logic of the Kripke frame $0<1$. The first and the last
case are known. New are the third case, the corresponding result for
all propositional theories --- the translation commutes with the
imaginaries exactly when no three prime filters form a chain ---, and
the fact that a non-invertible
endomorphism excludes commutation for every logic other than the
classical one.
\end{abstract}

\maketitle

\section{Introduction}

For an intermediate logic $L$ and a small coherent category $\cC$ let
$\HH_L\cC$ be the syntactic category of the internal coherent theory of
$\cC$ in first-order logic over $L$, and $\PP$ the pretopos
completion, which adds the imaginaries. There is a full and faithful
comparison $c^L_\cC\colon\HH_L\PP\cC\to\PP\HH_L\cC$ \cite{Interm}; if
$\cC$ is a pretopos, it is the unit $\HH_L\cC\to\PP\HH_L\cC$, and it is
an equivalence if and only if $\HH_L\cC$ is a pretopos
\cite[Lemma~2.2]{Interm}. We say that \emph{the translation into $L$
commutes with the imaginaries of $\cC$} if $c^L_\cC$ is an equivalence.

Let $\cX$ be a finite category in which idempotents split and
$\cP=\Fin^{\cX}$ the pretopos of the functors from $\cX$ to finite
sets; its models are, up to isomorphism, the evaluations at the
objects of $\cX$, with which we identify them, and the homomorphisms
of models are the arrows \cite[\S7]{Heyting}. An arrow is
\emph{invertible} if it is an isomorphism; $\cX$ is an
\emph{EI-category} if every endomorphism is invertible.

\begin{thm}\label{thm:class}
Let $\cX$ be a finite category in which idempotents split, $\cP=
\Fin^{\cX}$, and $L$ an intermediate logic. Then $c^L_\cP$ is an
equivalence if and only if one of the following holds:
\begin{enumerate}
\item[(1)] $L$ is classical logic;
\item[(2)] $L$ does not prove the weak excluded middle $\neg p\vee\neg
\neg p$, and $\cX$ is a groupoid;
\item[(3)] $L$ proves the weak excluded middle and not the linearity
axiom $(p\to q)\vee(q\to p)$, and $\cX$ is an EI-category without two
composable non-invertible arrows;
\item[(4)] $L$ proves linearity and is not classical, $\cX$ is an
EI-category, and no leg tree of $\cX$ is an obstruction for $n=n_L$ in
the sense of \cite[Thm.~6.15]{Strings}.
\end{enumerate}
\end{thm}

Here $n_L\in\{2,3,\dots,\infty\}$ is the supremum of the numbers of
elements of the finite chains which are frames of $L$. The condition in (4) is empty
if $n_L=2$ and holds if $\cX$ has no two composable non-invertible
arrows \cite[Cor.~6.3]{Strings}; so for an EI-category without two
composable non-invertible arrows the comparison is an equivalence for
every logic proving the weak excluded middle, and the answers in the
cases (3) and (4) can differ only when there are composable
non-invertible arrows.

We recall the condition in (4). A \emph{string} is a sequence of
non-invertible arrows $t_0\to t_1\to\dots\to t_r$ with at most $n$
objects, and $\cS_n$ is the category of strings, one in each
isomorphism class, an arrow $t\to s$ being an arrow $t_0\to s_0$ of
$\cX$ induced by an isomorphism of a final segment of $t$ onto $s$.
Fix a string $s$; for a string $t$ with $K_t:=\cS_n(t,s)\ne\emptyset$
let $V_t=\cX(t_0,s_0)$. A \emph{leg tree} $T$ at $t$ is a finite rooted
tree whose non-root nodes are labelled by arrows $\beta\colon t'\to
t''$ of $\cS_n$ into the string of the parent, $t'$ having one object
more than $t''$; it defines equivalence relations $E_T$ on $V_t$ and
$F_T$ on $K_t$ (and so do the subtrees $T'$ at the strings $t'$ of the
children of the root, on $V_{t'}$ and $K_{t'}$), both the equality for a tree without children and
otherwise the joins, in the lattices of equivalence relations on
$V_t$ and on $K_t$, over the children $(\beta,T')$ of the root, of
the relations $g\beta\mathrel{E_{T'}}g'\beta$, respectively
$g\beta\mathrel{F_{T'}}g'\beta$; and $T$ is an \emph{obstruction} if
$F_T$ is not the restriction of $E_T$ to $K_t$. The condition in (4)
is that this happens for no $s$, $t$ and $T$.
Here w.e.m.\ is the weak excluded middle.
\begin{center}
\begin{tabular}{lll}
$L$&a frame of $L$&equivalence iff\\\hline
classical&the point&always\\
no w.e.m.&the fork&$\cX$ is a groupoid\\
w.e.m., not linear&the diamond&$\cX$ is EI, no two composable\\
&&non-invertible arrows\\
linear, not classical&chains, $\le n_L$ elements&$\cX$ is EI, no
obstruction
\end{tabular}
\end{center}

\subsection*{What is new and what is known}
Known: case (1) \cite[Prop.~5.3]{Strings}; case (2), for every coherent
category in place of $\cP$, $c^L$ being an equivalence exactly in the
Boolean case \cite[Thm.~1.1(1)]{Interm}; case (4) for EI-categories
\cite[Thm.~1.1, Thm.~6.15]{Strings}; the fact that a non-invertible
endomorphism excludes an equivalence for the logics proving linearity
\cite[Thm.~7.1]{Strings}; and the obstruction over the diamond
\cite[Prop.~5.1]{Interm}, which was applied there to the propositional
theory of the four-element chain. New, to the author's knowledge:
Theorem~\ref{thm:nonEI}, that a non-invertible endomorphism excludes
an equivalence for every non-classical $L$, by the argument of
\cite[Thm.~7.1]{Strings}; and case (3), that is
Theorem~\ref{thm:wem}: the description of the models of $\HH_L\cP$ for
the logics proving the weak excluded middle when no two non-invertible
arrows are composable (Lemmas~\ref{lem:classical}
and~\ref{lem:top}), and the resulting equivalence. For a finite poset
$\cX$ the pretopos $\cP$ is the pretopos completion of a finite
propositional theory, that of the distributive lattice of the up-sets
of $\cX$, and case (3)
says that the comparison is an equivalence exactly when $\cX$ has no
chain of three elements; this answers, for finite propositional
theories and for the theories $\Fin^{\cX}$, the question of
\cite[\S8]{Interm} which non-Boolean theories have an equivalence in
this case, where the three-element chain as a distributive lattice,
whose poset of prime filters is the chain of two elements, was left
undecided. The same holds for all propositional theories:
Theorem~\ref{thm:propositional}, also new, says that in case (3) the
comparison for a bounded distributive lattice is an equivalence
exactly when there is no chain of three prime filters.

\section{The setting}\label{sec:setting}

We use the conventions and the results of \cite[\S2]{Strings}, which
hold for every intermediate logic $L$. The objects of $\HH_L\cP$ are
the subobjects $U\le A$, $A$ in $\cP$, obtained from the subobjects in
$\cP$ by finite meets and joins, $\Rightarrow$, and $f^*$,
$\exists_f$, $\forall_f$ for the morphisms $f$ of $\cP$. A
\emph{model} is a coherent functor $\hat G\colon\HH_L\cP\to\Set$; its
restriction to $\cP$ is the evaluation at an object $m$, and $\hat
G(U)$ is a subset of $A(m)$. A \emph{homomorphism} $\hat G\to\hat G'$
of models over $m$ and $m'$ is an arrow $\varphi\colon m\to m'$ of
$\cX$ with $A(\varphi)(\hat G(U))\subseteq\hat G'(U)$ for all $U$. We
write $y(x)=\cX(x,-)$.

\begin{lem}[{\cite[Lemma~2.2]{Strings}}]\label{lem:truth}
Let $\hat G$ be a model over $m$, $f\colon B\to A$ in $\cP$, $U,V\le
B$, $\alpha\in A(m)$ and $\beta\in B(m)$. Then $\alpha\in\hat G(\forall_f
U)$ if and only if $\beta'\in\hat G'(U)$ for every homomorphism
$\varphi\colon\hat G\to\hat G'$ and every $\beta'\in B(m')$ with
$f(\beta')=A(\varphi)(\alpha)$; and $\beta\in\hat G(U\Rightarrow V)$ if
and only if for every homomorphism $\varphi\colon\hat G\to\hat G'$,
$B(\varphi)(\beta)\in\hat G'(U)$ implies $B(\varphi)(\beta)\in\hat
G'(V)$.
\end{lem}

The models jointly reflect the order of subobjects, by the
completeness theorem of coherent logic. If $K$ is a finite poset such
that every $\uparrow\!k$ is a frame of $L$ and $M\colon K\to\cX$ a
functor, the forcing relation of the Kripke model $M$ defines a model
of $\HH_L\cP$ at every node $k$, over $M_k$, and the transitions are
homomorphisms \cite[Lemma~2.3]{Interm}. For a maximal node this is the
\emph{classical model} $\mathrm{cl}_m$ over $m=M_k$, the Heyting
functor $\HH_L\cP\to\Set$ extending the evaluation at $m$; for an
invertible arrow $\psi\colon m\to m'$ we have $A(\psi)(\mathrm{cl}_m(U))
=\mathrm{cl}_{m'}(U)$. For an arrow $a\colon x\to z$ which is not
invertible, the chain $0<1$ with $M_0=x$, $M_1=z$ and the transition
$a$ is a Kripke model whenever $L$ is not classical, the two-element
chain being a frame of every intermediate logic other than classical
logic; we write $\hat G_a$ for its model at the node $0$, so that
$a\colon\hat G_a\to\mathrm{cl}_z$ is a homomorphism.

We recall the frames \cite[Lemma~3.1]{Interm}: the fork is a frame of
$L$ if and only if $L$ does not prove the weak excluded middle; if $L$
proves it and not linearity, the diamond is a frame of $L$.
\[
\begin{tikzcd}[column sep=small,row sep=small]
\bullet&&\bullet&&&&\bullet&\\
&\bullet\ar[ul]\ar[ur]&&&&\bullet\ar[ur]&&\bullet\ar[ul]\\
&&&&&&\bullet\ar[ul]\ar[ur]&\\[-1ex]
&\text{fork}&&&&&\text{diamond}&
\end{tikzcd}
\]

\section{A non-invertible endomorphism}

\begin{thm}\label{thm:nonEI}
Let $L$ be an intermediate logic other than classical logic. If $\cX$
is not an EI-category, then $c^L_\cP$ is not an equivalence.
\end{thm}
\begin{proof}
This is the proof of \cite[Thm.~7.1]{Strings}, where $L$ is assumed to
prove linearity; we repeat it to make clear that this is not used. Let
$f\colon m\to m$ be a non-invertible endomorphism. A power $e$ of $f$
is idempotent, and $e\ne\id$. Let $r\colon m\to k$, $i\colon k\to m$
with $ri=\id_k$ and $ir=e$; then $i$ is not invertible, and we have the
homomorphism $i\colon\hat G_i\to\mathrm{cl}_m$.

For $w\in\cX(m,k)$ let $\Gamma_w\le y(m)\times y(k)$ be the subobject in
$\cP$ of the pairs $(z,x)$ with $z=xw$, and
\[
Q:=\exists_{p}\bigwedge_{w\in\cX(m,k)}\neg\Gamma_w\ \le\ y(k),
\]
$p$ being the projection to $y(k)$. Then $i\in\mathrm{cl}_m(Q)$: for
$z=\id_m$ we have $\id_m\ne iw$ for all $w$, since otherwise $i$ would
be invertible. And $\id_k\notin\hat G_i(Q)$: for every $z\in\cX(m,k)$
the pair $(z,\id_k)$ lies in $\Gamma_z$.

Let $X=y(k)\sqcup y(k)$ with the inclusions $j_1,j_2$, and $E\le
X\times X$ the join of the diagonal and of the images of $Q$ under
$(j_1,j_2)$ and $(j_2,j_1)$, an equivalence relation in $\HH_L\cP$. If
$\HH_L\cP$ were a pretopos, the quotient $X/E$ would be a subobject of
some $A$ in $\cP$, and the map $\hat G_i(X/E)\to\mathrm{cl}_m(X/E)$
induced by the homomorphism $i$ would be a restriction of $A(i)$,
which is injective, as $A(r)A(i)=\id$. But models preserve quotients;
the classes of $j_1(\id_k)$ and $j_2(\id_k)$ are different in $\hat
G_i(X/E)$ and are sent to the classes of $j_1(i)$ and $j_2(i)$, which
are equal in $\mathrm{cl}_m(X/E)$.
\end{proof}

\section{The weak excluded middle}

In this section $\cX$ is an EI-category. Then a split monomorphism or
epimorphism is invertible, a composite is invertible only if its
factors are, and the non-invertible arrows out of $x$ form a subfunctor
$N_x\le y(x)$ \cite[\S4]{Strings}. For an arrow $a\colon x\to z$ of any category $\cX$ let
$I_a\le y(x)$ be the image of $y(a)\colon y(z)\to y(x)$, the arrows
which factor through $a$.

\begin{lem}\label{lem:classical}
Let $L$ be an intermediate logic, $\cX$ a finite category in which
idempotents split, and $\hat G$ a model of $\HH_L\cP$ over $m$ such
that every homomorphism out of $\hat G$ has an invertible underlying
arrow. Then $\hat G=\mathrm{cl}_m$.
\end{lem}
\begin{proof}
If $\psi\colon\hat G\to\hat G'$ is a homomorphism, then $\hat G'$ has
the same property, a homomorphism $\lambda$ out of $\hat G'$ giving
the homomorphism $\lambda\psi$ out of $\hat G$. We show $\hat G(U)=
\mathrm{cl}_m(U)$ for all models with this property at once, by
induction on the construction of $U$. This is clear for the subobjects
in $\cP$ and compatible with the coherent operations. By
Lemma~\ref{lem:truth}, $\alpha\in\hat G(\forall_fU)$ if and only if
$\beta'\in\hat G'(U)=\mathrm{cl}_{m'}(U)$ for all homomorphisms $\psi
\colon\hat G\to\hat G'$ and all $\beta'$ over $A(\psi)(\alpha)$; since
$\psi$ is invertible and the classical models are compatible with the
invertible arrows, this holds if and only if $\beta\in\mathrm{cl}_m(U)$
for all $\beta$ over $\alpha$, the identity being among the $\psi$;
that is, if and only if $\alpha\in\mathrm{cl}_m(\forall_fU)$. The case
of $\Rightarrow$ is the same.
\end{proof}

\begin{lem}\label{lem:top}
Let $L$ be an intermediate logic other than classical logic which
proves the weak excluded middle, and let $\cX$ be a finite
EI-category without two composable non-invertible arrows. Let $\hat G$
be a model of $\HH_L\cP$ over $x$, and $R$ the set of the
non-invertible arrows out of $x$ which are homomorphisms out of $\hat
G$. Then either $R=\emptyset$ and $\hat G=\mathrm{cl}_x$, or $R$ is the
set of the arrows $\iota a$, for one $a\colon x\to z$ and all
invertible $\iota$ out of $z$, and $\hat G=\hat G_a$.
\end{lem}
\begin{proof}
If $a\colon x\to z$ is not invertible, every arrow out of $z$ is
invertible; so by Lemma~\ref{lem:classical} the target of a
homomorphism $a\in R$ is $\mathrm{cl}_z$. If $R=\emptyset$, then
$\hat G=\mathrm{cl}_x$ by Lemma~\ref{lem:classical}.

Let $a\colon x\to z$ and $a'\colon x\to z'$ be in $R$. If $a'=\lambda
a$ for some $\lambda$, then $\lambda$ is invertible. Suppose that
there is no such $\lambda$; then $a'\notin I_a(z')$, so
$a'\in\mathrm{cl}_{z'}(\neg I_a)$, while $a\in\mathrm{cl}_z(I_a)$. In
$\HH_L\cP$ the subobject $\neg I_a\vee\neg\neg I_a$ of $y(x)$ is the
top element; so $\id_x\in\hat G(\neg I_a)$, and then $a\in
\mathrm{cl}_z(\neg I_a)$ by the homomorphism $a$, or $\id_x\in\hat
G(\neg\neg I_a)$, and then $a'\in\mathrm{cl}_{z'}(\neg\neg I_a)$ by
$a'$. Both are impossible. Conversely $\iota a\in R$ for $a\in R$ and
$\iota$ invertible, as $\iota\colon\mathrm{cl}_z\to\mathrm{cl}_{z'}$ is
a homomorphism. So $R$, if not empty, is the set of the $\iota a$.

Let $R\ne\emptyset$, $a\in R$, and let $\psi\colon\hat G\to\hat G'$ be
a homomorphism with $\psi\colon x\to x'$ invertible; let $R'$ be the
set defined for $\hat G'$ as $R$ for $\hat G$. Then $R'\psi\subseteq
R$. We show $R'\ne\emptyset$. Otherwise all homomorphisms out of $\hat
G'$ are invertible in $\cX$, so that $\psi\in\hat G'(\neg N_x)$ by
Lemma~\ref{lem:truth}, as an invertible arrow is not in $N_x$; and
$a\in\mathrm{cl}_z(N_x)$. As before $\neg N_x\vee\neg\neg N_x$ is the
top element; $\id_x\in\hat G(\neg N_x)$ contradicts the homomorphism
$a$, and $\id_x\in\hat G(\neg\neg N_x)$ the homomorphism $\psi$. Hence
$R'$ is the set of the $\iota a\psi^{-1}$.

Now we show $\hat G(U)=\hat G_a(U)$ for all models $\hat G$ with
$R\ne\emptyset$ and all $a\in R$ at once, by induction on the
construction of $U$; note that $\psi\colon\hat G_a\to\hat G_{a\psi^{-1}}$
is an isomorphism of models for $\psi$ invertible, the two Kripke
models being isomorphic. The homomorphisms out of $\hat G$ are those
with an invertible underlying arrow, $\psi\colon\hat G\to\hat G'$ with $\hat G'(U)=\hat
G_{a\psi^{-1}}(U)$ by the induction hypothesis, among them the
identity, and the arrows $\iota a$ into classical models. So by
Lemma~\ref{lem:truth}, $\alpha\in\hat G(\forall_fU)$ if and only if
$\beta\in\hat G_a(U)$ for all $\beta$ over $\alpha$ and $\beta'\in
\mathrm{cl}_z(U)$ for all $\beta'$ over $A(a)(\alpha)$, which is the
forcing clause for $\forall_fU$ at the node $0$ of the Kripke model of
$a$. The case of $\Rightarrow$ is the same.
\end{proof}

\begin{thm}\label{thm:wem}
Let $L$ be an intermediate logic proving the weak excluded middle,
other than classical logic, and $\cX$ a finite EI-category.
\begin{enumerate}
\item[(a)] If $\cX$ has no two composable non-invertible arrows, then
$c^L_\cP$ is an equivalence.
\item[(b)] If $L$ does not prove linearity and $\cX$ has two composable
non-invertible arrows, then $c^L_\cP$ is not an equivalence.
\end{enumerate}
\end{thm}
\begin{proof}
(a) Let $L'$ be the logic of the Kripke frame $0<1$; it proves linearity,
$n_{L'}=2$, and $L\subseteq L'$. By Lemma~\ref{lem:top}, applied to $L$
and to $L'$, the models of $\HH_L\cP$ and of $\HH_{L'}\cP$ are the same
families of subsets: the classical models and the models $\hat G_a$,
whose values are defined by forcing and do not depend on the logic.
The axioms of $L'$ are valid in all these models; as the models of
$\HH_L\cP$ jointly reflect the order of subobjects, the axioms of $L'$
hold in $\HH_L\cP$, so that the canonical functor $\HH_L\cP\to
\HH_{L'}\cP$ is an equivalence. By
\cite[Cor.~6.3]{Strings} $c^{L'}_\cP$ is an equivalence.

(b) Let $f\colon x\to y$ and $g\colon y\to z$ be non-invertible. The
diamond is a frame of $L$. We apply \cite[Prop.~5.1]{Interm} to
$A=y(x)$, the subobjects $S=I_f$ and $T=I_{gf}$, the homomorphisms
$x\xrightarrow{f}y\xrightarrow{g}z$ of models of $\cP$, and the element
$\id_x\in A(x)$: we have $\id_x\notin I_f(x)$, as $f$ is not a split
monomorphism; $f\in I_f(y)$; $f\notin I_{gf}(y)$, since $f=hgf$ would
make $g$ a split monomorphism, the endomorphism $hg$ of $y$ being
invertible; and $gf\in I_{gf}(z)$. So $c^L_\cP$ is not an equivalence.
\end{proof}

The Kripke model over the diamond used in \cite[Prop.~5.1]{Interm} is
\[
\begin{tikzcd}[column sep=small,row sep=small]
&z&\\
y\ar[ur,"g"]&&x\ar[ul,"gf"']\\
&x\ar[ul,"f"]\ar[ur,"\id"']&
\end{tikzcd}
\]
and the proof of (b) is an instance of the principle behind
Theorem~\ref{thm:nonEI}: the transition $\id_x$ from the root to the
node on the right is a homomorphism of models whose underlying arrow
is a split monomorphism, and it does not reflect the value of
$I_f\Rightarrow I_{gf}$ at $\id_x$, which holds at the node on the
right and fails at the root because of the node on the left. For the
logics proving linearity the diamond is not a frame, and the first
obstructions appear with strings of three objects and non-trivial
automorphisms \cite[\S6]{Strings}.

\section{Propositional theories}\label{sec:prop}

The same argument, with prime filters in place of models, settles
case (3) for all propositional coherent theories. Let $D$ be a bounded
distributive lattice, regarded as a coherent category, $G_L$ the free
$L$-algebra on the bounded lattice $D$, with $\eta\colon D\to G_L$; as
a preorder $\HH_LD$ is $G_L$ \cite[\S6]{Interm}. For a prime filter $F$
of $G_L$ the filter $\eta^{-1}F$ of $D$ is prime; we say that $F$ is
\emph{over} $\eta^{-1}F$. We use the prime filter theorem in the form:
$u\to v\in F$ if and only if $u\in F'$ implies $v\in F'$ for every
prime filter $F'\supseteq F$ of $G_L$. A construction of free
Heyting algebras on distributive lattices by Esakia duality, adapted
to the algebras of the weak excluded middle and of linearity, is
given in \cite{Almeida}; we do not use it.

For a prime filter $P$ of $D$ let $\mathrm{cl}_P\colon G_L\to2$ be the
homomorphism of Heyting algebras extending the characteristic map of
$P$, and for prime filters $P\subsetneq P'$ let $h_{P,P'}\colon G_L\to
3$ be the homomorphism to the three-element chain $0<\frac12<1$, the
algebra of the up-sets of the Kripke frame $0<1$, which extends the
map sending $d$ to $1$ if $d\in P$, to $\frac12$ if $d\in P'\setminus
P$ and to $0$ otherwise; these exist if $L$ is not classical, the
three-element chain being then an $L$-algebra.

\begin{thm}\label{thm:propositional}
Let $L$ be an intermediate logic which proves the weak excluded middle
and is not classical, and $D$ a bounded distributive lattice.
\begin{enumerate}
\item[(a)] If $D$ has no chain $P_0\subsetneq P_1\subsetneq P_2$ of
prime filters, then $c^L_D$ is an equivalence.
\item[(b)] If $L$ does not prove linearity and $D$ has such a chain,
then $c^L_D$ is not an equivalence.
\end{enumerate}
\end{thm}
\begin{proof}
(b) Let $s\in P_1\setminus P_0$ and $t\in P_2\setminus P_1$. Then
\cite[Prop.~5.1]{Interm} applies to the terminal object $A$, the
subobjects $S=s$ and $T=t$, the homomorphisms $P_0\subseteq P_1\subseteq
P_2$ of models of $D$ and the unique element of $A(P_0)$.

(a) Every element of $G_L$ is obtained from elements $\eta(d)$ by
$\wedge$, $\vee$ and $\to$. We show that every prime filter $F$ of
$G_L$ is $\mathrm{cl}_P^{-1}(1)$ or $h_{P,P'}^{-1}(1)$.

\emph{Step 1.} If all prime filters $F'\supseteq F$ are over the same
$P$, then $F=\mathrm{cl}_P^{-1}(1)$. The prime filters above such an
$F$ have the same property; we show $u\in F$ iff $\mathrm{cl}_P(u)=1$
for all of them at once, by induction on $u$. For $u\to v$ this
follows from the prime filter theorem and the induction hypothesis for
the $F'\supseteq F$.

\emph{Step 2.} Let $F$ be over $P$ and $R$ the set of the prime filters
$P'\ne P$ of $D$ such that some prime filter $F'\supseteq F$ is over
$P'$. A filter $P'\in R$ contains $P$ strictly, so it is maximal among
the prime filters of $D$, by the hypothesis, and every prime filter
over $P'$ is $\mathrm{cl}_{P'}^{-1}(1)$ by Step 1. If $R=\emptyset$,
then $F=\mathrm{cl}_P^{-1}(1)$ by Step 1.

$R$ has at most one element: let $P',P''\in R$ be different, with
$F',F''\supseteq F$ over them, and $d\in P'\setminus P''$. As
$\neg\eta(d)\vee\neg\neg\eta(d)=1$, the prime filter $F$ contains
$\neg\eta(d)$, which is impossible as $\eta(d)\in F'$, or
$\neg\neg\eta(d)$, which is impossible as $\mathrm{cl}_{P''}(\neg\neg
\eta(d))=0$.

Let $R=\{P'\}$, and let $F^\flat\supseteq F$ be a prime filter over
$P$. Then the set $R^\flat$ defined for $F^\flat$ is $\{P'\}$: it is
contained in $R$, and if it were empty, then $F^\flat=\mathrm{cl}_P^{-1}
(1)$; for $d\in P'\setminus P$ the filter $F$ would contain $\neg\eta
(d)$, impossible as before, or $\neg\neg\eta(d)$, impossible as
$\mathrm{cl}_P(\neg\neg\eta(d))=0$.

\emph{Step 3.} If $R=\{P'\}$, then $F=h_{P,P'}^{-1}(1)$. We show $u\in
F$ iff $h_{P,P'}(u)=1$ for all prime filters $F$ with non-empty $R$ at
once, by induction on $u$. The prime filters above $F$ are those over
$P$, for which the induction hypothesis holds by Step 2, and those
over $P'$, which are $\mathrm{cl}_{P'}^{-1}(1)$. So $u\to v\in F$ iff
$h_{P,P'}(u)=1$ implies $h_{P,P'}(v)=1$, and $\mathrm{cl}_{P'}(u)=1$
implies $\mathrm{cl}_{P'}(v)=1$; as $\mathrm{cl}_{P'}(w)=1$ iff
$h_{P,P'}(w)\ge\frac12$, the map $3\to2$ sending $\frac12$ to $1$
being a homomorphism of Heyting algebras, this says $h_{P,P'}(u)\le h_{P,P'}(v)$, that
is $h_{P,P'}(u\to v)=1$.

Let $L'$ be the logic of the Kripke frame $0<1$; then $L\subseteq L'$,
and $2$ and $3$ are $L'$-algebras. If $u\not\le v$ in $G_L$, a prime
filter contains $u$ and not $v$, so the images of $u$ and $v$ under a
homomorphism to $2$ or $3$ which factors through $G_{L'}$ are
different. Hence the canonical homomorphism $G_L\to G_{L'}$, which is
surjective as both algebras are generated by the image of $D$, is an
isomorphism, and $\HH_LD\to\HH_{L'}D$ is an equivalence. Now $c^{L'}_D$
is an equivalence since $L'$ proves linearity
\cite[Thm.~6.3]{Interm}. The canonical functor $q\colon\HH_L\PP D\to
\HH_{L'}\PP D$ is the identity on formulas, hence essentially
surjective, and $c^{L'}_Dq$ is isomorphic to $c^L_D$ followed by the
equivalence $\PP\HH_LD\to\PP\HH_{L'}D$; so $c^L_D$ is essentially
surjective, and it is full and faithful \cite[Lemma~2.2]{Interm}.
\end{proof}

So for the logics of case (3) the translation commutes with the
imaginaries of a propositional theory if and only if the prime
spectrum of its Lindenbaum algebra has no chain of three points; for
the logics proving linearity it commutes for all propositional
theories, and for the logics of case (2) only for the Boolean ones
\cite{Interm}.

\section{The classification}

\begin{lem}\label{lem:boolean}
$\cP=\Fin^{\cX}$ is Boolean if and only if $\cX$ is a groupoid.
\end{lem}
\begin{proof}
If $\cX$ is a groupoid, the complement of a subfunctor is a subfunctor.
Let $a\colon x\to z$ be a non-invertible arrow. If $a$ is a split
monomorphism, $ra=\id$, then $ar$ is a non-invertible idempotent,
which is not a split monomorphism; so we may assume that $a$ is not a
split monomorphism. Then $\id_x\notin I_a(x)$, and a subfunctor of
$y(x)$ which contains $\id_x$ contains $a$; so $I_a$ has no complement.
\end{proof}

\begin{proof}[Proof of Theorem~\ref{thm:class}]
For classical logic $\HH_L\cP$ is a pretopos \cite[Prop.~5.3]{Strings}.
If $L$ does not prove the weak excluded middle, $c^L_\cP$ is an
equivalence if and only if $\cP$ is Boolean
\cite[Thm.~1.1(1)]{Interm}, that is, by Lemma~\ref{lem:boolean}, if and
only if $\cX$ is a groupoid. Let $L$ prove the weak excluded middle and
be different from classical logic. If $\cX$ is not an EI-category,
$c^L_\cP$ is not an equivalence by Theorem~\ref{thm:nonEI}. If $\cX$ is
an EI-category and $L$ does not prove linearity, Theorem~\ref{thm:wem}
gives (3). If $L$ proves linearity, $c^L_\cP$ is an equivalence if and
only if no leg tree is an obstruction \cite[Thm.~1.1,
Thm.~6.15]{Strings}.
\end{proof}

\begin{ex}\label{ex:small}
Small categories which separate the four cases of
Theorem~\ref{thm:class}, the columns giving the answer for the logics
of each case; in the last line
$\cX$ has two objects $m$, $k$ and the arrows $\id_m$, $e=ir$,
$r\colon m\to k$, $i\colon k\to m$, $\id_k=ri$:
\begin{center}
\begin{tabular}{l@{\ \ }c@{\ \ }c@{\ \ }c@{\ \ }c}
$\cX$&(1)&(2)&(3)&(4)\\\hline
a finite groupoid&yes&yes&yes&yes\\
the chain $0<1$&yes&no&yes&yes\\
the chain $0<1<2$&yes&no&no&yes\\
a split idempotent $e\ne\id$&yes&no&no&no
\end{tabular}
\end{center}
The entries for the chains in the last column are those of
\cite[Thm.~1.1(3)]{Interm}, for the chain $0<1<2$ in the third column
that of \cite[Thm.~1.1(2)]{Interm}, where it appears as the
four-element chain, the distributive lattice of its up-sets; the entry
for $0<1$ in the third column is new, and the last line is
Theorem~\ref{thm:nonEI}. In the last column the logic is not
classical; the answer depends on $n_L$ only for categories with
composable non-invertible arrows and non-trivial automorphisms, such
as \cite[Ex.~6.8]{Strings}, where it is yes for $n_L=2$ and no for
$n_L\ge3$; these have a no in the third column.
\end{ex}

So among the theories $\Fin^{\cX}$ every non-classical logic proving
linearity commutes with the imaginaries of a larger class than the
logics which do not prove it, in accordance with the propositional
case \cite{Interm}, this class depending on $n_L$; and for each of
the three types of non-classical logics the class is described by the
shape of the strings of non-invertible arrows: none, single arrows,
or strings subject to a condition on automorphisms.

\section{Questions}

\begin{que}
In case (3), for which coherent categories $\cC$ which are not of the
form $\Fin^{\cX}$ is $c^L_\cC$ an equivalence? For the theory of pure
equality it is not \cite[Thm.~7.2]{Interm}. Is the condition that the
homomorphisms of models of $\HH_L\cC$ whose underlying homomorphisms
of models of $\cC$ have a left inverse in the category of models of
$\cC$ reflect all first-order formulas over $L$, which
is necessary by the proofs of Theorems~\ref{thm:nonEI}
and~\ref{thm:wem}(b), also sufficient in some generality?
\end{que}

\begin{que}
Describe the models of $\HH_L\cP$ when $\cX$ is not an EI-category.
There are models with non-invertible endomorphisms
\cite[Ex.~7.2]{Strings}.
\end{que}

\section*{Declarations}
The author has no competing interests to declare.

\begin{thebibliography}{9}
\bibitem{Almeida} R.~N.~Almeida, \emph{Colimits of Heyting algebras
through Esakia duality}, preprint, arXiv:2402.08058.
\bibitem{Heyting} R.~Shibusawa, \emph{The free Heyting category on a
pretopos is a pretopos only in the Boolean case: the intuitionistic
translation and the pretopos completion}, preprint (2026), archived at
DOI 10.5281/zenodo.24182122.
\bibitem{Interm} R.~Shibusawa, \emph{Intermediate logics and
imaginaries: the translation into an intermediate logic commutes with
the pretopos completion of all finite propositional theories exactly
for the extensions of G\"odel--Dummett logic}, preprint (2026),
archived at DOI 10.5281/zenodo.24150610.
\bibitem{Strings} R.~Shibusawa, \emph{G\"odel--Dummett logic over a
finite category of models: the models of the translated theory are
strings of arrows, and a subgroup obstruction to commuting with the
imaginaries}, preprint (2026), archived at DOI
10.5281/zenodo.24327653.
\end{thebibliography}

\end{document}
